% chapters/12/section-01.tex
\section{Pekerjaan yang Berpindah Tangan}
\label{sec:perpindahan-tangan}

Bagian~2.6 telah menjelaskan perbedaan cara kerja pustaka dan mesin permainan
secara umum. Bab ini memperlihatkannya secara nyata, pada perkakas yang akan
Anda pakai sampai akhir semester.

\subsection{Apa yang Tidak Lagi Anda Tulis}

Bandingkan apa yang Anda tulis pada Bab~IV sampai Bab~VII dengan apa yang
tersedia pada Unity.

\begin{table}[htbp]
\centering
\caption{Pekerjaan yang pada Pygame Anda tulis sendiri, dan padanannya pada Unity}
\label{tab:pekerjaan-berpindah}
{\setlength{\tabcolsep}{4pt}\footnotesize%
\begin{tabularx}{\textwidth}{@{}L{4.2cm} L{3.2cm} Y@{}}
\toprule
\headrow \textbf{Pekerjaan} & \textbf{Pada Pygame} & \textbf{Pada Unity} \\
\midrule
Siklus permainan & Ditulis sendiri (Bab~IV) &
Dijalankan mesin; kode Anda dipanggil pada tahap yang telah ditentukan. \\
\addlinespace
Perhitungan fisika & Ditulis sendiri (Bab~VI) &
Dihitung mesin melalui komponen \istilah{Rigidbody}. \\
\addlinespace
Deteksi tabrakan & Ditulis sendiri (Bab~VII) &
Diperiksa mesin melalui \istilah{Collider}, dan dilaporkan kepada kode Anda
sebagai peristiwa. \\
\addlinespace
Penayangan & Dipanggil sendiri setiap bingkai &
Dikerjakan mesin; kode Anda hanya mengubah keadaan benda. \\
\addlinespace
Penyuntingan level & Diubah melalui kode &
Diubah melalui penyunting visual, tanpa menulis kode. \\
\bottomrule
\end{tabularx}}
\end{table}

Perpindahan tangan itu membawa dua akibat yang berlawanan. Di satu sisi,
pekerjaan Anda menjadi jauh lebih singkat: benda yang pada Bab~VII menuntut
puluhan baris kode kini cukup disusun dengan beberapa kali klik. Di sisi lain,
Anda kehilangan sebagian kejelasan: ketika perilaku permainan tidak sesuai
harapan, penyebabnya dapat berada pada kode Anda, pada pengaturan komponen,
atau pada cara mesin bekerja.

Pengetahuan yang Anda peroleh pada blok Python itulah bekal untuk membedakan
ketiganya. Kebiasaan mengukur, memasang pemeriksa mandiri, dan mencatat
keadaan---yang telah dibangun sejak Bab~IV---justru menjadi lebih penting pada
blok ini, sebab mesin menyembunyikan lebih banyak daripada yang ditampilkannya.

\begin{catatabox}
Satu hal tidak berpindah tangan: \textbf{kecepatan dinyatakan
dalam satuan per detik, dan dikalikan selang waktu}. Pada Unity, selang waktu
itu disediakan mesin melalui \code{Time.deltaTime}, bukan diukur sendiri
melalui \code{clock.tick()} seperti pada Bab~V. Konsepnya sama persis; yang
berubah hanyalah siapa yang mengukurnya.
\end{catatabox}

\subsection{Perkakas yang Diperlukan}

Unity dipasang melalui dua program yang terpisah, dan perbedaan keduanya perlu
diketahui sejak awal.

\begin{description}
  \item[Unity Hub] Program kecil yang mengurus pemasangan dan pengelolaan
        versi Unity beserta \istilah{project} Anda. Melalui program inilah Anda
        memasang versi Unity, membuka \istilah{project}, dan berpindah di antara
        beberapa \istilah{project}.
  \item[Unity Editor] Program utama tempat Anda bekerja: menyusun scene,
        menempelkan komponen, menulis \istilah{script}, dan menjalankan
        permainannya.
\end{description}

Versi yang dipakai sebaiknya versi \istilah{LTS} (\istilah{long-term support}),
yaitu versi yang didukung dalam jangka waktu panjang. Versi itu lebih stabil
dan lebih banyak dipakai, sehingga contoh kode dan tutorial yang Anda temukan
di internet lebih mungkin cocok.

\begin{catatabox}
Unity menuntut ruang penyimpanan dan memori yang jauh lebih besar daripada
Python dan Pygame: sekitar lima sampai lima belas gigabita untuk pemasangan,
bergantung pada modul yang dipilih. Periksa spesifikasi komputer Anda lebih
dahulu, dan pilih hanya modul yang memang diperlukan, misalnya dukungan
\istilah{2D} dan penyunting kode.
\end{catatabox}
